% Options for packages loaded elsewhere
\PassOptionsToPackage{unicode}{hyperref}
\PassOptionsToPackage{hyphens}{url}
\documentclass[aspectratio=169,ignorenonframetext]{beamer}
\usepackage{ma2003b-beamer}
\hypersetup{
  pdftitle={Multivariate Regression},
  pdfauthor={Juliho Castillo Colmenares},
  hidelinks,
  pdfcreator={LaTeX via pandoc}}

\title{\texorpdfstring{Multivariate
Regression}{Multivariate Regression}}
\author{\texorpdfstring{Juliho Castillo
Colmenares}{Juliho Castillo Colmenares}}
\subtitle{Advanced Methods for Multivariate Analysis}
\institute{Tec de Monterrey}
\date{}

\begin{document}
\frame{\titlepage}

\begin{frame}{Today's Agenda}
\begin{enumerate}
\item
  Logistic Regression Model
\item
  Inferences for Variances and Covariance Matrices
\item
  Inferences for a Vector of Means
\item
  MANOVA (Multivariate Analysis of Variance)
\item
  Canonical Correlation Analysis
\item
  Factor Analysis with Regression
\item
  Programming and Commercial Systems
\end{enumerate}
\end{frame}

\begin{frame}{Case Study: Healthcare Risk Assessment}
\textbf{Companion Example Throughout This Presentation}

\textbf{Synthetic data:} All patients and measurements below are
simulated for teaching purposes. Results illustrate the statistical
methods, not real clinical findings.

A simulated hospital scenario analyzing cardiovascular disease (CVD)
risk using multivariate methods to:

\begin{itemize}
\item
  Predict high-risk patients from lifestyle and physiological factors
\item
  Compare health profiles between risk groups
\item
  Evaluate the effectiveness of lifestyle interventions
\item
  Understand complex relationships between multiple health variables
\end{itemize}

Synthetic dataset with 1,000 simulated patients demonstrating practical
applications of each technique
\end{frame}

\begin{frame}{The Research Question}
\textbf{Scenario (synthetic data):} Hospital evaluating cardiovascular
disease (CVD) risk

\textbf{Dataset:}

\begin{itemize}
\item
  1,000 patients
\item
  13 predictor variables
\item
  Multiple health outcomes
\item
  Treatment intervention study
\end{itemize}
\end{frame}

\begin{frame}{Variables in Our Study}
\textbf{Lifestyle Factors:}

\begin{itemize}
\item
  \textbf{Age:} Years lived, affects baseline health risk
\item
  \textbf{BMI:} Body Mass Index, ratio of weight to height indicating
  body composition
\item
  \textbf{Exercise hours/week:} Physical activity level, linked to
  cardiovascular health
\item
  \textbf{Smoking years:} Duration of tobacco use, major risk factor for
  multiple diseases
\item
  \textbf{Alcohol consumption:} Frequency/amount of alcohol intake,
  impacts liver and heart health
\item
  \textbf{Stress score:} Quantified psychological stress level, affects
  blood pressure and hormones
\item
  \textbf{Sleep hours:} Nightly sleep duration, influences metabolic and
  cardiovascular function
\end{itemize}

\textbf{Physiological Measurements:}

\begin{itemize}
\item
  \textbf{Blood pressure (systolic/diastolic):} Force of blood against
  artery walls (peak/resting)
\item
  \textbf{Cholesterol:} Blood lipid levels, excess increases arterial
  plaque buildup
\item
  \textbf{Glucose:} Blood sugar level, indicator of diabetes risk and
  metabolic function
\item
  \textbf{Triglycerides:} Fat molecules in blood, high levels increase
  heart disease risk
\item
  \textbf{HDL:} ``Good cholesterol'' that removes harmful lipids from
  bloodstream
\end{itemize}
\end{frame}

\begin{frame}{Research Objectives}
\begin{enumerate}
\item
  \textbf{Predict} CVD risk from patient characteristics
\item
  \textbf{Compare} health profiles between risk groups
\item
  \textbf{Evaluate} lifestyle intervention effectiveness
\item
  \textbf{Understand} relationships between lifestyle and physiology
\item
  \textbf{Validate} assumptions for multivariate tests
\end{enumerate}
\end{frame}

\begin{frame}{Logistic Regression}
Moving Beyond Linear Regression

\textbf{What is Logistic Regression?}

A statistical method for modeling binary outcomes (Yes/No,
Success/Failure, 0/1) using predictor variables

\textbf{Key Features:}

\begin{itemize}
\item
  Predicts probabilities (0 to 1) rather than continuous values
\item
  Uses the logistic (sigmoid) function to model the probability of an
  event
\item
  Estimates coefficients via maximum likelihood (not least squares)
\item
  Widely used in classification problems: medical diagnosis, credit
  scoring, marketing
\end{itemize}

\textbf{When to Use:} When your outcome variable is categorical
(especially binary)
\end{frame}

\begin{frame}{When Linear Regression Fails}
\textbf{Problem:} Binary outcomes (Yes/No, Success/Failure, 0/1)

Using ordinary least squares (OLS) for binary responses creates
fundamental problems:

\textbf{Prediction Issues:}

\begin{itemize}
\item
  Predicted probabilities can be negative or exceed 1 (nonsensical
  values)
\item
  No mechanism to constrain predictions to {[}0, 1{]}
\end{itemize}

\textbf{Assumption Violations:}

\begin{itemize}
\item
  Response is not continuous (violates normality assumption)
\item
  Errors follow Bernoulli distribution, not Normal distribution
\item
  Heteroscedastic errors (variance depends on X:
  \(\text{Var}\left( Y|X \right) = p(X)\left( 1 - p(X) \right)\))
\item
  Non-constant variance violates homoscedasticity assumption
\end{itemize}

Linear regression is mathematically possible but statistically
inappropriate for binary data.
\end{frame}

\begin{frame}{Logistic Regression Solution}
Our goal is to find a function that maps predictor values to
probabilities while staying within {[}0, 1{]}. Linear regression fails
here because predictions can fall outside this valid probability range.

Given a binary outcome variable \(Y\) (taking values 0 or 1 for failure
or success) and a vector of predictor variables
\(X = \left( X_{1},X_{2},\ldots,X_{p} \right)\), we model the
probability of success as:

\[P\left( Y = 1~|~X \right) = p(X)\]

where \(0 \leq p(X) \leq 1\) represents the probability that \(Y = 1\)
given the values of \(X\). This approach ensures our predictions always
remain valid probabilities.
\end{frame}

\begin{frame}{The Logit Transformation}
To connect probabilities (bounded between 0 and 1) to a linear
combination of predictors (unbounded), we use the \textbf{logit
transformation}, also known as the \textbf{log-odds}.

Recall that \(p(X) = P\left( Y = 1~|~X \right)\) is our probability of
success. The odds of success are \(\frac{p(X)}{1 - p(X)}\), and taking
the natural logarithm gives us:

\[\text{ logit}\left( p(X) \right) = \log(\frac{p(X)}{1 - p(X)}) = \beta_{0} + \beta_{1}X_{1} + \ldots + \beta_{p}X_{p}\]

where:

\begin{itemize}
\item
  \(p(X)\) is the same probability function defined in the previous
  slide
\item
  \(\beta_{0},\beta_{1},\ldots,\beta_{p}\) are coefficients to be
  estimated
\item
  The logit maps probabilities from {[}0, 1{]} to
  \(( - \infty, + \infty)\)
\item
  This transformation makes the model linear in the parameters
\end{itemize}
\end{frame}

\begin{frame}{The Logit Transformation}
\textbf{Properties:}

\begin{itemize}
\item
  Maps {[}0,1{]} to \(( - \infty, + \infty)\)
\item
  Linear in parameters
\item
  Interpretable as the log-odds (a coefficient's exponential is the odds
  \textbf{ratio} per unit increase)
\end{itemize}
\end{frame}

\begin{frame}{The Logistic Function}
\textbf{Inverse Logit:} Solving for \(p(X)\) from the logit equation
gives the logistic function.

\textbf{Notation:} \(X = \left( X_{1},X_{2},\ldots,X_{p} \right)\) is
the vector of \(p\) predictor variables. We write \(p(X)\) to mean the
probability depends on all predictors.

Define the linear combination:
\(\eta = \beta_{0} + \beta_{1}X_{1} + \beta_{2}X_{2} + \ldots + \beta_{p}X_{p}\)

Then the logistic function is:

\[p(X) = \frac{e^{\eta}}{1 + e^{\eta}} = \frac{1}{1 + e^{- \eta}}\]

This is the \textbf{logistic} or \textbf{sigmoid} function, guaranteeing
\(p(X) \in \lbrack 0,1\rbrack\) for any predictor values.
\end{frame}

\begin{frame}{Maximum Likelihood Estimation}
Since we cannot use ordinary least squares, we estimate coefficients
using \textbf{Maximum Likelihood Estimation (MLE)}.

\textbf{Our Data:} We have collected \(n\) observations. Each
observation consists of:

\begin{itemize}
\item
  An outcome: \(y_{i} \in \left\{ 0,1 \right\}\) (e.g., \(y_{i} = 1\)
  means ``has disease'', \(y_{i} = 0\) means ``no disease'')
\item
  Predictor values: \(x_{i1},x_{i2},\ldots,x_{ip}\) (the values of \(p\)
  predictors for person \(i\))
\end{itemize}

\textbf{Unknown Parameters:} We need to estimate \((p + 1)\)
coefficients

\begin{itemize}
\item
  \(\beta_{0}\) = intercept coefficient
\item
  \(\beta_{1},\beta_{2},\ldots,\beta_{p}\) = coefficients for the \(p\)
  predictors
\end{itemize}

We write
\(\mathbf{\beta} = \left( \beta_{0},\beta_{1},\ldots,\beta_{p} \right)^{T}\)
for the entire vector of unknown parameters.
\end{frame}

\begin{frame}{Maximum Likelihood Estimation}
\textbf{The Probability Model:} For each observation \(i\), the
probability of success depends on:

\begin{enumerate}
\item
  The predictor values for that observation:
  \(x_{i1},x_{i2},\ldots,x_{ip}\)
\item
  The unknown parameters: \(\beta_{0},\beta_{1},\ldots,\beta_{p}\)
\end{enumerate}

Define the linear combination:
\(\eta_{i} = \beta_{0} + \beta_{1}x_{i1} + \beta_{2}x_{i2} + \ldots + \beta_{p}x_{ip}\)

The probability is computed using the logistic function:

\[P\left( Y_{i} = 1~|~x_{i1},\ldots,x_{ip};\beta_{0},\ldots,\beta_{p} \right) = \frac{1}{1 + e^{- \eta_{i}}}\]

We write this as \(\pi_{i}\left( \mathbf{\beta} \right)\) for brevity,
emphasizing it depends on the parameters \(\mathbf{\beta}\).

\textbf{Important:} Each \(\pi_{i}\left( \mathbf{\beta} \right)\)
depends on BOTH the observed predictors for observation \(i\) AND the
unknown parameters \(\mathbf{\beta}\).
\end{frame}

\begin{frame}{Maximum Likelihood Estimation}
\textbf{Probability of Observing Outcome \(y_{i}\):} Given parameters
\(\mathbf{\beta}\), the probability of observing the actual outcome
\(y_{i}\) for observation \(i\) is formally written as
\(P\left( Y_{i} = y_{i}~|~\ldots \right)\), where \(Y_{i}\) is the
random variable for the outcome. The following is a common shorthand:

\[P\left( Y_{i} = y_{i}~|~\mathbf{x}_{i},\mathbf{\beta} \right) = \left\lbrack \pi_{i}\left( \mathbf{\beta} \right) \right\rbrack^{y_{i}}\left\lbrack 1 - \pi_{i}\left( \mathbf{\beta} \right) \right\rbrack^{1 - y_{i}}\]

This formula evaluates to:

\begin{itemize}
\item
  \(\pi_{i}\left( \mathbf{\beta} \right)\) when the observed outcome is
  \(y_{i} = 1\) (success)
\item
  \(1 - \pi_{i}\left( \mathbf{\beta} \right)\) when the observed outcome
  is \(y_{i} = 0\) (failure)
\end{itemize}

\textbf{Key Idea of MLE:} Find the parameter values \(\mathbf{\beta}\)
that make the observed data \(\left( y_{1},y_{2},\ldots,y_{n} \right)\)
most probable.
\end{frame}

\begin{frame}{Maximum Likelihood Estimation}
\textbf{Likelihood Function:} Assuming observations are independent, the
probability of observing ALL our data is:

\[L\left( \mathbf{\beta} \right) = \prod_{i = 1}^{n}P\left( Y_{i} = y_{i}~|~\mathbf{x}_{i},\mathbf{\beta} \right) = \prod_{i = 1}^{n}\left\lbrack \pi_{i}\left( \mathbf{\beta} \right) \right\rbrack^{y_{i}}\left\lbrack 1 - \pi_{i}\left( \mathbf{\beta} \right) \right\rbrack^{1 - y_{i}}\]

\textbf{Log-Likelihood Function:} Taking natural logarithm (easier to
maximize):

\[\ell(\mathbf{\beta}) = \sum_{i = 1}^{n}\left\lbrack y_{i}\log(\pi_{i}\left( \mathbf{\beta} \right)) + \left( 1 - y_{i} \right)\log(1 - \pi_{i}\left( \mathbf{\beta} \right)) \right\rbrack\]

\textbf{Goal:} Find \(\mathbf{\beta}^{\ast}\) that maximizes
\(\ell(\mathbf{\beta})\). This requires numerical optimization
(Newton-Raphson, gradient descent, etc.) since no closed-form solution
exists.
\end{frame}

\begin{frame}{Interpreting Coefficients: Log-Odds vs. Odds Ratios}
Recall the logit equation:
\[\log(\frac{p}{1 - p}) = \beta_{0} + \beta_{1}X_{1} + \ldots + \beta_{p}X_{p}\]

\textbf{1. The Raw Coefficient (\(\beta_{j}\))}

\begin{itemize}
\item
  A one-unit increase in \(X_{j}\) corresponds to a \(\beta_{j}\) change
  in the \textbf{log-odds} of the outcome.
\item
  This is mathematically precise but not intuitive. A ``0.2 change in
  log-odds'' is hard to grasp.
\end{itemize}

\textbf{2. The Odds Ratio (\(e^{\beta_{j}}\)) -- A Better Way}

\begin{itemize}
\item
  By exponentiating the coefficient, we get the \textbf{Odds Ratio}
  (OR).
\item
  The OR tells us the \textbf{multiplicative change} in the
  \textbf{odds} of the outcome for a one-unit increase in \(X_{j}\).
\end{itemize}

\textbf{Example: Stress Score (\(\beta_{\text{stress }} = 0.223\))}\\
The odds ratio is \(e^{0.223} \approx 1.25\).\\
\textbf{Interpretation:} For each one-point increase in the stress
score, the \textbf{odds} of having high CVD risk are multiplied by 1.25
(i.e., increase by 25\%), holding other variables constant.
\end{frame}

\begin{frame}{Interpreting Coefficients: The Odds Ratio (cont.)}
\textbf{Interpreting Different Odds Ratios:}

{\def\LTcaptype{none} % do not increment counter
\begin{longtable}[]{@{}
  >{\centering\arraybackslash}p{(\linewidth - 4\tabcolsep) * \real{0.2501}}
  >{\centering\arraybackslash}p{(\linewidth - 4\tabcolsep) * \real{0.2501}}
  >{\raggedright\arraybackslash}p{(\linewidth - 4\tabcolsep) * \real{0.5001}}@{}}
\toprule\noalign{}
\endhead
\textbf{If \(\beta_{j}\) is\ldots{}} & \textbf{Then \(e^{\beta_{j}}\)
is\ldots{}} & \textbf{Meaning for a 1-unit increase in
\(X_{j}\)\ldots{}} \\
Positive (\(> 0\)) & Greater than 1 & The odds of the outcome
\textbf{increase}. (e.g., OR=1.25 means 25\% increase in odds) \\
Negative (\(< 0\)) & Less than 1 & The odds of the outcome
\textbf{decrease}. (e.g., OR=0.80 means 20\% decrease in odds) \\
Zero & Exactly 1 & There is no change in the odds of the outcome. \\
\bottomrule\noalign{}
\end{longtable}
}

\textbf{Example: Exercise (\(\beta_{\text{exercise }} = - 0.328\))}\\
The odds ratio is \(e^{- 0.328} \approx 0.72\).\\
\textbf{Interpretation:} For each additional hour of exercise per week,
the \textbf{odds} of having high CVD risk are multiplied by 0.72 (i.e.,
decrease by 28\%), holding other variables constant.
\end{frame}

\begin{frame}{Model Fit and Diagnostics}
Assessing how well our logistic regression model explains the data is
crucial. Unlike linear regression, we use different metrics.

\textbf{Deviance:}

\begin{itemize}
\item
  A key metric for models where ordinary least squares doesn't apply
  (like logistic regression). It's a measure of how much your model's
  predictions deviate from the actual data.
\item
  Think of it as the equivalent of the \textbf{Sum of Squared Errors}
  you see in linear regression.
\item
  A lower deviance means your model fits the data more closely.
\item
  For a deeper dive into deviance, see this video:
  \href{https://www.youtube.com/watch?v=JC56jS2gVUE}{Deviance Residuals
  Explained}
\end{itemize}

\textbf{Pseudo R-squared:}

\begin{itemize}
\item
  While linear regression has a standard R-squared to measure
  goodness-of-fit, logistic and other generalized linear models use
  ``Pseudo'' R-squared.
\item
  It tells you how much better your model is than a ``null'' model that
  only includes an intercept (essentially, a model that just guesses the
  average outcome).
\item
  \textbf{McFadden's \(R^{2}\)} is a popular version:
  \[R_{\text{McF }}^{2} = 1 - \frac{\log(\mathcal{L}_{\text{model}})}{\log(\mathcal{L}_{\text{null}})}\].
\item
  A higher value (closer to 1) indicates a better fit, but the scale is
  not directly comparable to the R-squared from linear regression.
\end{itemize}
\end{frame}

\begin{frame}{Classification Performance}
\textbf{Confusion Matrix:}

{\def\LTcaptype{none} % do not increment counter
\begin{longtable}[]{@{}
  >{\centering\arraybackslash}p{(\linewidth - 4\tabcolsep) * \real{0.3333}}
  >{\centering\arraybackslash}p{(\linewidth - 4\tabcolsep) * \real{0.3333}}
  >{\centering\arraybackslash}p{(\linewidth - 4\tabcolsep) * \real{0.3333}}@{}}
\toprule\noalign{}
\endhead
& \textbf{Predicted 0} & \textbf{Predicted 1} \\
\textbf{Actual 0} & True Negative (TN) & False Positive (FP) \\
\textbf{Actual 1} & False Negative (FN) & True Positive (TP) \\
\bottomrule\noalign{}
\end{longtable}
}
\end{frame}

\begin{frame}{Classification Metrics}
\textbf{Accuracy:} \(\frac{\text{TP } + \text{ TN}}{n}\)

\textbf{Sensitivity (Recall):}
\(\frac{\text{TP }}{\text{TP } + \text{ FN}}\)

\textbf{Specificity:} \(\frac{\text{TN }}{\text{TN } + \text{ FP}}\)

\textbf{Precision:} \(\frac{\text{TP }}{\text{TP } + \text{ FP}}\)
\end{frame}

\begin{frame}{Case Study: Predicting CVD Risk}
\textbf{Application of Logistic Regression}

Objective: Predict high CVD risk (0/1) from 13 patient characteristics
\end{frame}

\begin{frame}{CVD Prediction: Model Setup}
\textbf{Predictors (13 variables):}

\begin{itemize}
\item
  Demographics: age, BMI
\item
  Lifestyle: exercise, smoking, alcohol, stress, sleep
\item
  Physiology: BP, cholesterol, glucose, triglycerides, HDL
\end{itemize}

\textbf{Outcome:} CVD risk high (binary: 0 = low risk, 1 = high risk)

\textbf{Data split:} 70\% training (n=700), 30\% testing (n=300)
\end{frame}

\begin{frame}{Top Risk Factors: Odds Ratios}
{\def\LTcaptype{none} % do not increment counter
\begin{longtable}[]{@{}
  >{\raggedright\arraybackslash}p{(\linewidth - 4\tabcolsep) * \real{0.4444}}
  >{\centering\arraybackslash}p{(\linewidth - 4\tabcolsep) * \real{0.2222}}
  >{\raggedright\arraybackslash}p{(\linewidth - 4\tabcolsep) * \real{0.3333}}@{}}
\toprule\noalign{}
\endhead
\textbf{Predictor} & \textbf{Odds Ratio} & \textbf{Interpretation} \\
Exercise hours/week & 0.72 & 28\% lower odds per hour \\
Stress score & 1.25 & 25\% higher odds per point \\
Sleep hours & 0.80 & 20\% lower odds per hour \\
BMI & 1.19 & 19\% higher odds per unit \\
\bottomrule\noalign{}
\end{longtable}
}

\textbf{All significant predictors contribute to risk assessment}
\end{frame}

\begin{frame}{CVD Prediction: Model Performance}
\textbf{Confusion Matrix (Test Set):}

{\def\LTcaptype{none} % do not increment counter
\begin{longtable}[]{@{}
  >{\centering\arraybackslash}p{(\linewidth - 4\tabcolsep) * \real{0.4288}}
  >{\centering\arraybackslash}p{(\linewidth - 4\tabcolsep) * \real{0.2859}}
  >{\centering\arraybackslash}p{(\linewidth - 4\tabcolsep) * \real{0.2859}}@{}}
\toprule\noalign{}
\endhead
& \textbf{Pred Low} & \textbf{Pred High} \\
\textbf{Actual Low} & 106 & 44 \\
\textbf{Actual High} & 42 & 108 \\
\bottomrule\noalign{}
\end{longtable}
}

\textbf{Metrics:}

\begin{itemize}
\item
  Accuracy: 71\%
\item
  AUC-ROC: 0.77
\item
  Balanced precision/recall
\end{itemize}
\end{frame}

\begin{frame}{Key Insights: CVD Prediction}
\begin{enumerate}
\item
  Exercise has the strongest \textbf{standardized} association with risk
  (OR = 0.72 per hour/week; strongest by per-SD effect size, not raw OR)
\item
  Stress also significantly increases risk (OR = 1.25 per point)
\item
  Raw odds ratios are not comparable across predictors on different
  scales -- ranking requires standardized effects
\item
  Model achieves good discrimination (AUC = 0.77) on this simulated
  dataset
\end{enumerate}

\textbf{Illustrative only:} This is a simulated dataset -- read as a
demonstration of the method, not a clinical finding
\end{frame}

\begin{frame}{Inferences for Covariance Matrices}
Testing Variability Structure
\end{frame}

\begin{frame}{Why Test Covariance Matrices?}
\textbf{Applications:}

\begin{itemize}
\item
  Homogeneity assumptions in MANOVA
\item
  Comparing variability between groups
\item
  Validating models
\item
  Quality control
\end{itemize}
\end{frame}

\begin{frame}{The Wishart Distribution}
\textbf{Multivariate Generalization of Chi-Square}

If \(X_{1},\ldots,X_{n} \sim N_{p(\mu,\Sigma)}\), then:

\[S \sim W_{p(n - 1,\Sigma)}\]

where
\(S = \sum_{i = 1}^{n}\left( X_{i} - \overline{X} \right)\left( X_{i} - \overline{X} \right)^{T}\)
\end{frame}

\begin{frame}{Testing Single Covariance Matrix}
\textbf{Null Hypothesis:}

\[H_{0}:\Sigma = \Sigma_{0}\]

\textbf{Test Statistic:} Based on likelihood ratio

\[\Lambda = |S\frac{|}{|}\Sigma_{0}|\]
\end{frame}

\begin{frame}{Box's M Test}
\textbf{Testing Equality of Covariance Matrices}

\[H_{0}:\Sigma_{1} = \Sigma_{2} = \ldots = \Sigma_{g}\]
\end{frame}

\begin{frame}{Box's M Test Statistic}
\[M = (n - g)\log|S_{\text{pooled}}| - \sum_{i = 1}^{g}\left( n_{i} - 1 \right)\log|S_{i}|\]

where:

\begin{itemize}
\item
  \(S_{i}\) = covariance matrix for group \(i\)
\item
  \(S_{\text{pooled}}\) = pooled covariance matrix
\end{itemize}
\end{frame}

\begin{frame}{Box's M Test Properties}
\textbf{Asymptotic Distribution:} The raw \(M\) must first be rescaled
by a correction factor \(c_{1}\); only the corrected
\(\chi^{2} = M\left( 1 - c_{1} \right)\) is chi-square distributed.
There is no valid raw-\(M\) cutoff.

\textbf{Limitation:} Very sensitive to normality violations

\textbf{Alternatives:} Permutation tests, robust methods
\end{frame}

\begin{frame}{Bartlett's Test for Univariate Data}
\textbf{Special Case:} Testing equality of variances (p=1)

\[H_{0}:\sigma_{1}^{2} = \sigma_{2}^{2} = \ldots = \sigma_{g}^{2}\]

\textbf{Test Statistic:} Chi-square distributed
\end{frame}

\begin{frame}{Case Study: Validating MANOVA Assumptions}
\textbf{Application of Box's M Test}

Question: Are covariance matrices equal between treatment groups (MANOVA
assumption)?
\end{frame}

\begin{frame}{Box's M Test: Setup}
\textbf{Testing Homogeneity of Covariances:}

\[H_{0}:\mathbf{\Sigma}_{\text{Control }} = \mathbf{\Sigma}_{\text{Intervention }}\]

\textbf{Variables (p = 4):}

\begin{itemize}
\item
  Systolic BP, Diastolic BP
\item
  Cholesterol, Glucose
\end{itemize}

\textbf{Groups:}

\begin{itemize}
\item
  Control: n = 479
\item
  Intervention: n = 521
\end{itemize}

\textbf{Why test?} MANOVA assumes equal covariance matrices across
groups
\end{frame}

\begin{frame}{Box's M Test: Results}
\textbf{Raw Statistic:} \(M = 8.49\) (not itself chi-square distributed)

\textbf{Corrected:} \(\chi^{2} = 8.46\), \(\text{df } = 10\),
\(p = 0.584\)

\textbf{Conclusion:} Fail to reject \(H_{0}\) (\(p = 0.584\)). No
evidence against equal covariance matrices.

\textbf{Note:} A raw-statistic cutoff like ``\(M < 30\)'' is not a valid
decision rule -- \(M\) must be corrected and compared to its chi-square
distribution. Failing to reject does not \textbf{prove} equality.
\end{frame}

\begin{frame}{Key Insights: Assumption Testing}
\begin{enumerate}
\item
  Box's M requires a correction factor before its chi-square
  approximation is valid
\item
  Corrected test found no evidence against equal covariances
  (\(p = 0.584\))
\item
  ``Fail to reject'' is not the same as ``proven equal''
\item
  Treatment groups show no detected difference in variability patterns
\end{enumerate}

\textbf{Methodological importance:} Always check assumptions using the
correct test statistic, and report results as ``fail to reject,'' not
``validated''
\end{frame}

\begin{frame}{Inferences for a Vector of Means}
Multivariate Hypothesis Testing
\end{frame}

\begin{frame}{From t-test to Hotelling's T-squared}
\textbf{Univariate:} t-test for single mean

\[t = \frac{\overline{x} - \mu_{0}}{\frac{s}{\sqrt{n}}}\]

\textbf{Multivariate:} Hotelling's \(T^{2}\) for mean vector
\end{frame}

\begin{frame}{Hotelling's T-squared Test}
\textbf{One-Sample Test:}

\[H_{0}:\mathbf{\mu} = \mathbf{\mu}_{0}\]

\textbf{Test Statistic:}

\[T^{2} = {n\left( \overline{\mathbf{X}} - \mathbf{\mu}_{0} \right)}^{T}S^{- 1}\left( \overline{\mathbf{X}} - \mathbf{\mu}_{0} \right)\]
\end{frame}

\begin{frame}{Distribution of T-squared}
\textbf{Transform to F Distribution:}

\[F = \frac{(n - p)T^{2}}{(n - 1)p} \sim F_{p,n - p}\]

where:

\begin{itemize}
\item
  \(p\) = number of variables
\item
  \(n\) = sample size
\end{itemize}
\end{frame}

\begin{frame}{Two-Sample Hotelling's T-squared}
\textbf{Testing Difference Between Groups:}

\[H_{0}:\mathbf{\mu}_{1} = \mathbf{\mu}_{2}\]
\end{frame}

\begin{frame}{Two-Sample T-squared Statistic}
\[T^{2} = \left( \frac{n_{1}n_{2}}{n_{1} + n_{2}} \right)\left( {\overline{\mathbf{X}}}_{1} - {\overline{\mathbf{X}}}_{2} \right)^{T}S_{\text{pooled}}^{- 1}\left( {\overline{\mathbf{X}}}_{1} - {\overline{\mathbf{X}}}_{2} \right)\]

\textbf{F Transformation:}

\[F = \frac{\left( n_{1} + n_{2} - p - 1 \right)T^{2}}{\left( n_{1} + n_{2} - 2 \right)p} \sim F_{p,n_{1} + n_{2} - p - 1}\]
\end{frame}

\begin{frame}{Confidence Region for Mean Vector}
\textbf{Multivariate Confidence Region:}

Ellipsoid centered at \(\overline{\mathbf{X}}\)

\[{n\left( \mathbf{\mu} - \overline{\mathbf{X}} \right)}^{T}S^{- 1}\left( \mathbf{\mu} - \overline{\mathbf{X}} \right) \leq \frac{(n - 1)p}{n - p}F_{\alpha;p,n - p}\]
\end{frame}

\begin{frame}{Simultaneous Confidence Intervals}
\textbf{Bonferroni Correction:}

For \(p\) variables, use \(\frac{\alpha}{p}\) for each interval

\textbf{T-squared Intervals:} More efficient but wider than individual
intervals
\end{frame}

\begin{frame}{Case Study: Comparing Risk Groups}
\textbf{Application of Hotelling's T-squared}

Question: Do high-risk and low-risk CVD patients differ in their
multivariate health profile?

\textbf{Caveat:} The risk-group label is itself computed from a score
containing systolic BP, cholesterol, and glucose -- three of the six
variables tested below. A significant result here demonstrates the test
recovering that construction rule, not an independently discovered
relationship.
\end{frame}

\begin{frame}{Health Profile Comparison: Setup}
\textbf{Two Groups:}

\begin{itemize}
\item
  Low risk: n = 500
\item
  High risk: n = 500
\end{itemize}

\textbf{Variables (p = 6):}

\begin{itemize}
\item
  Systolic BP, Diastolic BP
\item
  Cholesterol, Glucose
\item
  Triglycerides, HDL
\end{itemize}

\textbf{Goal:} Single omnibus test for all 6 variables simultaneously
\end{frame}

\begin{frame}{Mean Differences by Risk Group}
{\def\LTcaptype{none} % do not increment counter
\begin{longtable}[]{@{}
  >{\raggedright\arraybackslash}p{(\linewidth - 6\tabcolsep) * \real{0.4000}}
  >{\centering\arraybackslash}p{(\linewidth - 6\tabcolsep) * \real{0.2000}}
  >{\centering\arraybackslash}p{(\linewidth - 6\tabcolsep) * \real{0.2000}}
  >{\centering\arraybackslash}p{(\linewidth - 6\tabcolsep) * \real{0.2000}}@{}}
\toprule\noalign{}
\endhead
\textbf{Variable} & \textbf{Low Risk} & \textbf{High Risk} &
\textbf{Difference} \\
Systolic BP & 123.8 & 131.1 & +7.3 \\
Diastolic BP & 78.1 & 82.9 & +4.7 \\
Cholesterol & 184.1 & 196.4 & +12.3 \\
Glucose & 110.0 & 116.8 & +6.9 \\
Triglycerides & 133.6 & 145.5 & +11.9 \\
HDL & 44.5 & 41.2 & -3.3 \\
\bottomrule\noalign{}
\end{longtable}
}
\end{frame}

\begin{frame}{Hotelling's T-squared Results}
\textbf{Test Statistic:} \[T^{2} = 228.65\]

\textbf{F Transformation:} \[F = 37.92,\quad\text{ df } = (6,993)\]

\textbf{P-value:} \textless{} 0.0001

\textbf{Conclusion:} The test statistically confirms the health markers
differ between groups -- as expected, since the grouping rule was built
from three of them. This demonstrates the mechanics of Hotelling's
\(T^{2}\), not an independent clinical finding.
\end{frame}

\begin{frame}{Key Insights: Risk Group Differences}
\begin{enumerate}
\item
  High-risk patients show higher values across all adverse markers, by
  construction
\item
  Largest differences: cholesterol (+12.3 mg/dL) and triglycerides
  (+11.9 mg/dL)
\item
  HDL (protective) is lower in high-risk group (-3.3 mg/dL)
\item
  Multivariate test accounts for correlations among measurements
\end{enumerate}

\textbf{Reminder:} This recovers the simulated data-generating rule -- a
real analysis needs a grouping variable independent of the tested
variables
\end{frame}

\begin{frame}{MANOVA}
Multivariate Analysis of Variance
\end{frame}

\begin{frame}{What is MANOVA?}
\textbf{Extension of ANOVA to Multiple Dependent Variables}

\begin{itemize}
\item
  ANOVA: One response variable
\item
  MANOVA: Multiple response variables simultaneously
\end{itemize}
\end{frame}

\begin{frame}{Why Use MANOVA?}
\textbf{Instead of Multiple ANOVAs:}

\begin{enumerate}
\item
  Controls Type I error rate
\item
  Accounts for correlations among responses
\item
  More powerful when responses related
\item
  Tests overall group effect
\end{enumerate}
\end{frame}

\begin{frame}{MANOVA Model}
\textbf{One-Way MANOVA:}

\[\mathbf{Y}_{ij} = \mathbf{\mu} + \mathbf{\alpha}_{i} + \mathbf{\varepsilon}_{ij}\]

where:

\begin{itemize}
\item
  \(\mathbf{Y}_{ij}\) = response vector for observation \(j\) in group
  \(i\)
\item
  \(\mathbf{\mu}\) = overall mean vector
\item
  \(\mathbf{\alpha}_{i}\) = group effect vector
\item
  \(\mathbf{\varepsilon}_{ij}\) = error vector
\end{itemize}
\end{frame}

\begin{frame}{MANOVA Assumptions}
\begin{enumerate}
\item
  \textbf{Multivariate Normality:} Errors follow multivariate normal
\item
  \textbf{Independence:} Observations independent
\item
  \textbf{Homogeneity of Covariance:} Equal covariance matrices across
  groups
\end{enumerate}
\end{frame}

\begin{frame}{Testing Assumptions}
\textbf{Multivariate Normality:}

\begin{itemize}
\item
  Mardia's test
\item
  Q-Q plots for each variable
\end{itemize}

\textbf{Homogeneity:} Box's M test
\end{frame}

\begin{frame}{MANOVA Matrices}
\textbf{Between-Groups Matrix (H):}

\[\mathbf{H} = \sum_{i = 1}^{g}n_{i}\left( {\overline{\mathbf{Y}}}_{i} - \overline{\mathbf{Y}} \right)\left( {\overline{\mathbf{Y}}}_{i} - \overline{\mathbf{Y}} \right)^{T}\]

\textbf{Within-Groups Matrix (E):}

\[\mathbf{E} = \sum_{i = 1}^{g}\sum_{j = 1}^{n_{i}}\left( \mathbf{Y}_{ij} - {\overline{\mathbf{Y}}}_{i} \right)\left( \mathbf{Y}_{ij} - {\overline{\mathbf{Y}}}_{i} \right)^{T}\]
\end{frame}

\begin{frame}{Wilks' Lambda}
\textbf{Most Common Test Statistic:}

\[\Lambda = |\mathbf{E}\frac{|}{|}\mathbf{E} + \mathbf{H}|\]
\end{frame}

\begin{frame}{Wilks' Lambda Properties}
\textbf{Interpretation:}

\begin{itemize}
\item
  Range: {[}0, 1{]}
\item
  Small values: Strong group differences
\item
  Lambda = 1: No group differences
\end{itemize}

\textbf{Represents:} Proportion of total variance not explained by
groups
\end{frame}

\begin{frame}{Other MANOVA Test Statistics}
\textbf{Pillai's Trace:}

\[V = \text{ tr}\left( {\mathbf{H}(\mathbf{H} + \mathbf{E})}^{- 1} \right)\]

\textbf{Hotelling-Lawley Trace:}

\[U = \text{ tr}\left( \mathbf{H}\mathbf{E}^{- 1} \right)\]

\textbf{Roy's Largest Root:} Largest eigenvalue of
\(\mathbf{H}\mathbf{E}^{- 1}\)
\end{frame}

\begin{frame}{Choosing Test Statistic}
{\def\LTcaptype{none} % do not increment counter
\begin{longtable}[]{@{}
  >{\raggedright\arraybackslash}p{(\linewidth - 2\tabcolsep) * \real{0.4286}}
  >{\raggedright\arraybackslash}p{(\linewidth - 2\tabcolsep) * \real{0.5714}}@{}}
\toprule\noalign{}
\endhead
\textbf{Statistic} & \textbf{Best When} \\
Wilks' Lambda & General use (most common) \\
Pillai's Trace & Robust to violations \\
Hotelling-Lawley & Equal group sizes \\
Roy's Root & Group difference on one dimension \\
\bottomrule\noalign{}
\end{longtable}
}
\end{frame}

\begin{frame}{Post-Hoc Tests in MANOVA}
\textbf{After Significant MANOVA:}

\begin{enumerate}
\item
  Univariate ANOVAs (with correction)
\item
  Discriminant analysis
\item
  Contrast tests for specific hypotheses
\end{enumerate}
\end{frame}

\begin{frame}{Case Study: Treatment Intervention}
\textbf{Application of MANOVA}

Question: Does a lifestyle intervention improve multiple health outcomes
simultaneously?
\end{frame}

\begin{frame}{Treatment Intervention Study: Setup}
\textbf{Groups:}

\begin{itemize}
\item
  Control: n = 479 (standard care)
\item
  Intervention: n = 521 (lifestyle program)
\end{itemize}

\textbf{Outcomes (p = 4):}

\begin{itemize}
\item
  Systolic BP
\item
  Diastolic BP
\item
  Cholesterol
\item
  Glucose
\end{itemize}

\textbf{Why MANOVA?} Controls Type I error while testing all outcomes
together
\end{frame}

\begin{frame}{MANOVA Results: Treatment Effect}
\textbf{Wilks' Lambda:} \(\Lambda = 0.889\)

\textbf{F Statistic:} \(F = 31.05\), df = (4, 995)

\textbf{P-value:} \textless{} 0.0001

\textbf{Conclusion:} The intervention significantly improves health
outcomes across the multivariate profile

\textbf{Also significant:} Pillai's trace, Hotelling-Lawley, Roy's root
(all p \textless{} 0.0001)
\end{frame}

\begin{frame}{Mean Improvements by Treatment Group}
{\def\LTcaptype{none} % do not increment counter
\begin{longtable}[]{@{}
  >{\raggedright\arraybackslash}p{(\linewidth - 6\tabcolsep) * \real{0.3846}}
  >{\centering\arraybackslash}p{(\linewidth - 6\tabcolsep) * \real{0.1923}}
  >{\centering\arraybackslash}p{(\linewidth - 6\tabcolsep) * \real{0.1923}}
  >{\centering\arraybackslash}p{(\linewidth - 6\tabcolsep) * \real{0.2308}}@{}}
\toprule\noalign{}
\endhead
\textbf{Outcome} & \textbf{Control} & \textbf{Intervention} &
\textbf{Difference} \\
Systolic BP & 130.8 & 124.3 & -6.5\textsuperscript{***} \\
Diastolic BP & 82.7 & 78.5 & -4.2\textsuperscript{***} \\
Cholesterol & 194.4 & 186.5 & -7.9\textsuperscript{***} \\
Glucose & 116.3 & 110.7 & -5.6\textsuperscript{***} \\
\bottomrule\noalign{}
\end{longtable}
}

\textsuperscript{***} All differences significant at p \textless{} 0.001
in follow-up ANOVAs
\end{frame}

\begin{frame}{Key Insights: Intervention Effects}
\begin{enumerate}
\item
  In this simulation, intervention reduces all cardiovascular risk
  markers
\item
  Largest effect on blood pressure (-6.5 / -4.2 mmHg)
\item
  Differences also survive Bonferroni correction across the 4 follow-up
  ANOVAs
\item
  MANOVA provides single omnibus test (no Type I error inflation)
\end{enumerate}

\textbf{Note:} Synthetic data illustrating the method -- not a real
intervention-effectiveness finding
\end{frame}

\begin{frame}{Canonical Correlation Analysis}
Relating Two Sets of Variables
\end{frame}

\begin{frame}{What is Canonical Correlation?}
\textbf{Purpose:} Find maximum correlation between linear combinations
of two sets of variables

\begin{itemize}
\item
  Set 1: \(X_{1},X_{2},\ldots,X_{p}\)
\item
  Set 2: \(Y_{1},Y_{2},\ldots,Y_{q}\)
\end{itemize}
\end{frame}

\begin{frame}{Canonical Correlation vs. Other Methods}
{\def\LTcaptype{none} % do not increment counter
\begin{longtable}[]{@{}
  >{\raggedright\arraybackslash}p{(\linewidth - 4\tabcolsep) * \real{0.4288}}
  >{\centering\arraybackslash}p{(\linewidth - 4\tabcolsep) * \real{0.2859}}
  >{\centering\arraybackslash}p{(\linewidth - 4\tabcolsep) * \real{0.2859}}@{}}
\toprule\noalign{}
\endhead
\textbf{Method} & \textbf{Set 1} & \textbf{Set 2} \\
Correlation & 1 variable & 1 variable \\
Multiple Regression & Multiple & 1 variable \\
Canonical Correlation & Multiple & Multiple \\
\bottomrule\noalign{}
\end{longtable}
}
\end{frame}

\begin{frame}{Canonical Variates}
\textbf{First Canonical Variate Pair:}

\[U_{1} = a_{11}X_{1} + a_{12}X_{2} + \ldots + a_{1p}X_{p}\]
\[V_{1} = b_{11}Y_{1} + b_{12}Y_{2} + \ldots + b_{1q}Y_{q}\]

such that \(\text{cor}\left( U_{1},V_{1} \right)\) is maximized
\end{frame}

\begin{frame}{Number of Canonical Correlations}
\textbf{How Many Pairs?}

\[k = \min(p,q)\]

Each subsequent pair:

\begin{itemize}
\item
  Uncorrelated with previous pairs
\item
  Maximizes remaining correlation
\end{itemize}
\end{frame}

\begin{frame}{Canonical Correlation Coefficients}
\textbf{Ordering:}

\[\rho_{1} \geq \rho_{2} \geq \ldots \geq \rho_{k} \geq 0\]

where \(\rho_{i}\) is the \(i\)-th canonical correlation
\end{frame}

\begin{frame}{Testing Significance}
\textbf{Test All Correlations:}

\[H_{0}:\rho_{1} = \rho_{2} = \ldots = \rho_{k} = 0\]

\textbf{Test Remaining Correlations:}

\[H_{0}:\rho_{m + 1} = \ldots = \rho_{k} = 0\]
\end{frame}

\begin{frame}{Wilks' Lambda for Canonical Correlation}
\[\Lambda = \prod_{i = 1}^{k}\left( 1 - \rho_{i}^{2} \right)\]

Approximate chi-square distribution for testing
\end{frame}

\begin{frame}{Canonical Loadings}
\textbf{Structure Coefficients:}

Correlation between original variables and canonical variates

\begin{itemize}
\item
  Help interpret meaning of canonical variates
\item
  More stable than canonical weights
\end{itemize}
\end{frame}

\begin{frame}{Redundancy Analysis}
\textbf{Proportion of Variance Explained:}

How much variance in one set is explained by the other set through
canonical variates

\[\text{ Redundancy } = \left( \frac{1}{p} \right)\sum_{j = 1}^{p}R_{X_{j},V_{1}}^{2}\]
\end{frame}

\begin{frame}{Interpreting Canonical Correlations}
\begin{enumerate}
\item
  \textbf{Examine significance:} Are correlations statistically
  significant?
\item
  \textbf{Check magnitude:} Are correlations practically meaningful?
\item
  \textbf{Interpret loadings:} What do canonical variates represent?
\item
  \textbf{Assess redundancy:} How much variance explained?
\end{enumerate}
\end{frame}

\begin{frame}{Case Study: Lifestyle vs. Physiology}
\textbf{Application of Canonical Correlation}

Question: How do lifestyle factors relate to physiological health
markers?
\end{frame}

\begin{frame}{Lifestyle-Physiology Relationship: Setup}
\textbf{Set 1 - Lifestyle Factors (p = 5):}

\begin{itemize}
\item
  Exercise hours/week
\item
  Smoking years
\item
  Alcohol units/week
\item
  Stress score
\item
  Sleep hours
\end{itemize}

\textbf{Set 2 - Physiological Markers (q = 6):}

\begin{itemize}
\item
  Systolic BP, Diastolic BP
\item
  Cholesterol, Glucose
\item
  Triglycerides, HDL
\end{itemize}

\textbf{Maximum pairs:} min(5, 6) = 5
\end{frame}

\begin{frame}{Canonical Correlations: Results}
{\def\LTcaptype{none} % do not increment counter
\begin{longtable}[]{@{}
  >{\centering\arraybackslash}p{(\linewidth - 4\tabcolsep) * \real{0.3001}}
  >{\centering\arraybackslash}p{(\linewidth - 4\tabcolsep) * \real{0.3001}}
  >{\raggedright\arraybackslash}p{(\linewidth - 4\tabcolsep) * \real{0.4002}}@{}}
\toprule\noalign{}
\endhead
\textbf{Pair} & \textbf{Correlation} & \textbf{Interpretation} \\
1 & 0.639 & Strong relationship \\
2 & 0.244 & Moderate relationship \\
3-5 & \textless{} 0.12 & Weak relationships \\
\bottomrule\noalign{}
\end{longtable}
}

\textbf{Focus on first canonical correlation (r = 0.639)}
\end{frame}

\begin{frame}{First Canonical Variate: Lifestyle}
\textbf{Canonical Loadings (Structure Coefficients):}

{\def\LTcaptype{none} % do not increment counter
\begin{longtable}[]{@{}
  >{\raggedright\arraybackslash}p{(\linewidth - 2\tabcolsep) * \real{0.6667}}
  >{\centering\arraybackslash}p{(\linewidth - 2\tabcolsep) * \real{0.3333}}@{}}
\toprule\noalign{}
\endhead
\textbf{Variable} & \textbf{Loading} \\
Exercise hours & +0.65 \\
Stress score & -0.53 \\
Alcohol units & -0.37 \\
Sleep hours & +0.33 \\
Smoking years & -0.27 \\
\bottomrule\noalign{}
\end{longtable}
}

\textbf{Interpretation:} Healthy lifestyle pattern (more exercise, less
stress)
\end{frame}

\begin{frame}{First Canonical Variate: Physiology}
\textbf{Canonical Loadings:}

{\def\LTcaptype{none} % do not increment counter
\begin{longtable}[]{@{}
  >{\raggedright\arraybackslash}p{(\linewidth - 2\tabcolsep) * \real{0.6667}}
  >{\centering\arraybackslash}p{(\linewidth - 2\tabcolsep) * \real{0.3333}}@{}}
\toprule\noalign{}
\endhead
\textbf{Variable} & \textbf{Loading} \\
Diastolic BP & -0.70 \\
Systolic BP & -0.68 \\
Cholesterol & -0.65 \\
HDL & +0.65 \\
Glucose & -0.61 \\
Triglycerides & -0.59 \\
\bottomrule\noalign{}
\end{longtable}
}

\textbf{Interpretation:} Favorable health profile (lower BP, higher HDL)
\end{frame}

\begin{frame}{Key Insights: Lifestyle-Physiology Link}
\begin{enumerate}
\item
  Strong first-pair canonical correlation (r = 0.639) between lifestyle
  and health variates
\item
  Healthy lifestyle pattern → favorable physiological profile, in this
  simulation
\item
  Exercise and low stress load most heavily on the first lifestyle
  variate
\item
  Blood pressure and cholesterol load most heavily on the first
  physiology variate
\end{enumerate}

\textbf{Note:} Association between canonical variates, not a domain-wide
or causal claim
\end{frame}

\begin{frame}{Factor Analysis with Regression}
Combining Dimension Reduction and Prediction
\end{frame}

\begin{frame}{The Multicollinearity Problem}
\textbf{Issue:} Highly correlated predictors in regression

\textbf{Consequences:}

\begin{itemize}
\item
  Unstable coefficient estimates
\item
  Large standard errors
\item
  Difficult interpretation
\item
  Poor prediction in new samples
\end{itemize}
\end{frame}

\begin{frame}{Factor-Based Regression Solution}
\textbf{Strategy:}

\begin{enumerate}
\item
  Extract factors from correlated predictors
\item
  Use factor scores as predictors
\item
  Fit regression with orthogonal factors
\end{enumerate}
\end{frame}

\begin{frame}{Factor-Based Regression Workflow}
\begin{enumerate}
\item
  \textbf{Factor Analysis:} Extract factors from \(X\) variables
\item
  \textbf{Compute Factor Scores:} For each observation
\item
  \textbf{Regression:} Predict \(Y\) using factor scores
\item
  \textbf{Interpretation:} Results in terms of factors
\end{enumerate}
\end{frame}

\begin{frame}{Benefits of Factor-Based Regression}
\textbf{Advantages:}

\begin{itemize}
\item
  Reduces multicollinearity (orthogonal factors)
\item
  Dimensionality reduction (fewer predictors)
\item
  Conceptual interpretation (latent constructs)
\item
  More stable estimates
\end{itemize}
\end{frame}

\begin{frame}{Comparing Approaches}
{\def\LTcaptype{none} % do not increment counter
\begin{longtable}[]{@{}
  >{\raggedright\arraybackslash}p{(\linewidth - 4\tabcolsep) * \real{0.4288}}
  >{\centering\arraybackslash}p{(\linewidth - 4\tabcolsep) * \real{0.2859}}
  >{\centering\arraybackslash}p{(\linewidth - 4\tabcolsep) * \real{0.2859}}@{}}
\toprule\noalign{}
\endhead
\textbf{Aspect} & \textbf{Direct Regression} & \textbf{Factor
Regression} \\
Multicollinearity & Problem & Eliminated \\
Interpretation & Original variables & Latent factors \\
Predictors & Many & Few \\
Variance explained & Higher & May be lower \\
\bottomrule\noalign{}
\end{longtable}
}
\end{frame}

\begin{frame}{Principal Components Regression}
\textbf{Alternative Approach:}

Use PCA instead of factor analysis

\textbf{Difference:}

\begin{itemize}
\item
  PCA: Explains total variance
\item
  FA: Explains common variance (removes unique variance)
\end{itemize}
\end{frame}

\begin{frame}{Other Methods for Multicollinearity}
\textbf{Ridge Regression:} Shrinks coefficients toward zero

\textbf{Lasso:} Variable selection via L1 penalty

\textbf{Partial Least Squares:} Finds components that predict Y well
\end{frame}

\begin{frame}{Cautions and Limitations}
\textbf{Factor-Based Regression Limitations:}

\begin{itemize}
\item
  Factor extraction somewhat subjective
\item
  Results depend on specific sample
\item
  Prediction requires computing factor scores with same loadings
\item
  May lose some predictive information
\end{itemize}
\end{frame}

\begin{frame}{Programming and Commercial Systems}
Implementing Multivariate Methods
\end{frame}

\begin{frame}[fragile]{Python for Multivariate Analysis}
\textbf{Key Libraries:}

\begin{itemize}
\item
  \texttt{statsmodels}: Statistical models and tests
\item
  \texttt{scikit-learn}: Machine learning algorithms
\item
  \texttt{numpy} / \texttt{scipy}: Numerical computations
\item
  \texttt{pandas}: Data manipulation
\end{itemize}
\end{frame}

\begin{frame}[fragile]{Python: Logistic Regression}
\begin{Shaded}
\begin{Highlighting}[]
\ImportTok{from}\NormalTok{ sklearn.linear\_model }\ImportTok{import}\NormalTok{ LogisticRegression}

\NormalTok{model }\OperatorTok{=}\NormalTok{ LogisticRegression()}
\NormalTok{model.fit(X\_train, y\_train)}
\NormalTok{predictions }\OperatorTok{=}\NormalTok{ model.predict(X\_test)}
\NormalTok{probabilities }\OperatorTok{=}\NormalTok{ model.predict\_proba(X\_test)}
\end{Highlighting}
\end{Shaded}
\end{frame}

\begin{frame}[fragile]{Python: Hotelling's T-squared}
\begin{Shaded}
\begin{Highlighting}[]
\ImportTok{from}\NormalTok{ scipy.stats }\ImportTok{import}\NormalTok{ chi2}
\ImportTok{import}\NormalTok{ numpy }\ImportTok{as}\NormalTok{ np}

\CommentTok{\# Compute T{-}squared statistic}
\NormalTok{diff }\OperatorTok{=}\NormalTok{ mean1 }\OperatorTok{{-}}\NormalTok{ mean2}
\NormalTok{S\_pooled\_inv }\OperatorTok{=}\NormalTok{ np.linalg.inv(S\_pooled)}
\NormalTok{T2 }\OperatorTok{=}\NormalTok{ (n1 }\OperatorTok{*}\NormalTok{ n2) }\OperatorTok{/}\NormalTok{ (n1 }\OperatorTok{+}\NormalTok{ n2) }\OperatorTok{*}\NormalTok{ diff.T }\OperatorTok{@}\NormalTok{ S\_pooled\_inv }\OperatorTok{@}\NormalTok{ diff}

\CommentTok{\# Transform to F}
\NormalTok{p }\OperatorTok{=} \BuiltInTok{len}\NormalTok{(mean1)}
\NormalTok{F\_stat }\OperatorTok{=}\NormalTok{ ((n1 }\OperatorTok{+}\NormalTok{ n2 }\OperatorTok{{-}}\NormalTok{ p }\OperatorTok{{-}} \DecValTok{1}\NormalTok{) }\OperatorTok{*}\NormalTok{ T2) }\OperatorTok{/}\NormalTok{ ((n1 }\OperatorTok{+}\NormalTok{ n2 }\OperatorTok{{-}} \DecValTok{2}\NormalTok{) }\OperatorTok{*}\NormalTok{ p)}
\end{Highlighting}
\end{Shaded}
\end{frame}

\begin{frame}[fragile]{Python: MANOVA}
\begin{Shaded}
\begin{Highlighting}[]
\ImportTok{from}\NormalTok{ statsmodels.multivariate.manova }\ImportTok{import}\NormalTok{ MANOVA}

\CommentTok{\# Fit MANOVA model}
\NormalTok{manova }\OperatorTok{=}\NormalTok{ MANOVA.from\_formula(}
    \StringTok{\textquotesingle{}Y1 + Y2 + Y3 \textasciitilde{} Group\textquotesingle{}}\NormalTok{,}
\NormalTok{    data}\OperatorTok{=}\NormalTok{df}
\NormalTok{)}

\CommentTok{\# Test results}
\BuiltInTok{print}\NormalTok{(manova.mv\_test())}
\end{Highlighting}
\end{Shaded}
\end{frame}

\begin{frame}[fragile]{Python: Canonical Correlation}
\begin{Shaded}
\begin{Highlighting}[]
\ImportTok{from}\NormalTok{ sklearn.cross\_decomposition }\ImportTok{import}\NormalTok{ CCA}

\CommentTok{\# Canonical correlation analysis}
\NormalTok{cca }\OperatorTok{=}\NormalTok{ CCA(n\_components}\OperatorTok{=}\DecValTok{2}\NormalTok{)}
\NormalTok{cca.fit(X\_set, Y\_set)}

\CommentTok{\# Transform to canonical variates}
\NormalTok{X\_c, Y\_c }\OperatorTok{=}\NormalTok{ cca.transform(X\_set, Y\_set)}

\CommentTok{\# Canonical correlations}
\NormalTok{correlations }\OperatorTok{=}\NormalTok{ [np.corrcoef(X\_c[:, i], Y\_c[:, i])[}\DecValTok{0}\NormalTok{, }\DecValTok{1}\NormalTok{]}
                \ControlFlowTok{for}\NormalTok{ i }\KeywordTok{in} \BuiltInTok{range}\NormalTok{(}\DecValTok{2}\NormalTok{)]}
\end{Highlighting}
\end{Shaded}
\end{frame}

\begin{frame}[fragile]{R for Multivariate Analysis}
\textbf{Key Packages:}

\begin{itemize}
\item
  \texttt{stats}: Base statistical functions
\item
  \texttt{MASS}: Advanced statistical methods
\item
  \texttt{car}: Companion to Applied Regression
\item
  \texttt{vegan}: Multivariate analysis
\end{itemize}
\end{frame}

\begin{frame}[fragile]{R: MANOVA Example}
\begin{Shaded}
\begin{Highlighting}[]
\CommentTok{\# Fit MANOVA}
\NormalTok{model }\OtherTok{\textless{}{-}} \FunctionTok{manova}\NormalTok{(}\FunctionTok{cbind}\NormalTok{(Y1, Y2, Y3) }\SpecialCharTok{\textasciitilde{}}\NormalTok{ Group, }\AttributeTok{data =}\NormalTok{ df)}

\CommentTok{\# Test results}
\FunctionTok{summary}\NormalTok{(model, }\AttributeTok{test =} \StringTok{"Wilks"}\NormalTok{)}
\FunctionTok{summary}\NormalTok{(model, }\AttributeTok{test =} \StringTok{"Pillai"}\NormalTok{)}

\CommentTok{\# Follow{-}up univariate tests}
\FunctionTok{summary.aov}\NormalTok{(model)}
\end{Highlighting}
\end{Shaded}
\end{frame}

\begin{frame}{Commercial Software: SPSS}
\textbf{GUI-Based Analysis:}

\begin{itemize}
\item
  Analyze \textgreater{} General Linear Model \textgreater{}
  Multivariate
\item
  Analyze \textgreater{} Regression \textgreater{} Binary Logistic
\item
  Analyze \textgreater{} Correlate \textgreater{} Canonical Correlation
\end{itemize}

\textbf{Syntax:} Also supports command syntax for reproducibility
\end{frame}

\begin{frame}[fragile]{Commercial Software: SAS}
\textbf{Key Procedures:}

\begin{itemize}
\item
  \texttt{PROC\ LOGISTIC}: Logistic regression
\item
  \texttt{PROC\ GLM}: General linear models (MANOVA)
\item
  \texttt{PROC\ CANCORR}: Canonical correlation
\item
  \texttt{PROC\ FACTOR}: Factor analysis
\end{itemize}
\end{frame}

\begin{frame}{Software Comparison}
{\def\LTcaptype{none} % do not increment counter
\begin{longtable}[]{@{}
  >{\raggedright\arraybackslash}p{(\linewidth - 4\tabcolsep) * \real{0.2501}}
  >{\raggedright\arraybackslash}p{(\linewidth - 4\tabcolsep) * \real{0.3751}}
  >{\raggedright\arraybackslash}p{(\linewidth - 4\tabcolsep) * \real{0.3751}}@{}}
\toprule\noalign{}
\endhead
\textbf{Software} & \textbf{Strengths} & \textbf{Limitations} \\
Python & Free, flexible, ML integration & Statistical testing less
developed \\
R & Free, comprehensive stats & Steeper learning curve \\
SPSS & GUI, easy to learn & Expensive, less flexible \\
SAS & Enterprise, comprehensive & Very expensive, complex \\
\bottomrule\noalign{}
\end{longtable}
}
\end{frame}

\begin{frame}{Choosing Software}
\textbf{Considerations:}

\begin{itemize}
\item
  Cost (free vs. commercial)
\item
  Learning curve
\item
  Specific methods needed
\item
  Integration with workflow
\item
  Reproducibility requirements
\item
  Team expertise
\end{itemize}
\end{frame}

\begin{frame}{Best Practices: Code Documentation}
\textbf{Essential Elements:}

\begin{itemize}
\item
  Comment your code clearly
\item
  Document data preprocessing steps
\item
  Record package versions
\item
  Save random seeds for reproducibility
\item
  Version control (Git)
\end{itemize}
\end{frame}

\begin{frame}{Best Practices: Workflow}
\begin{enumerate}
\item
  \textbf{Data Cleaning:} Handle missing values, outliers
\item
  \textbf{Exploratory Analysis:} Visualize distributions
\item
  \textbf{Check Assumptions:} Test before analysis
\item
  \textbf{Run Analysis:} Use appropriate methods
\item
  \textbf{Validate Results:} Cross-validation, diagnostics
\item
  \textbf{Document:} Clear reporting
\end{enumerate}
\end{frame}

\begin{frame}{Case Study Summary}
\textbf{Healthcare Risk Assessment: What We Learned}
\end{frame}

\begin{frame}{Key Findings: Prediction and Classification}
\textbf{Logistic Regression Results:}

\begin{itemize}
\item
  71\% accuracy predicting CVD risk (AUC = 0.77) on synthetic data
\item
  Exercise strongest by standardized effect size (not raw OR)
\item
  Stress increases risk 25\% per point (OR = 1.25)
\item
  Demonstrates how such a model could identify high-risk patients on
  real data
\end{itemize}

\textbf{Note:} Simulated dataset -- illustrates the method, not a
clinical recommendation
\end{frame}

\begin{frame}{Key Findings: Group Comparisons}
\textbf{Hotelling's T-squared:}

\begin{itemize}
\item
  Groups differ significantly across 6 health markers (T² = 228.65, p
  \textless{} 0.0001)
\item
  Expected: 3 of these 6 variables define the grouping rule (circular by
  construction)
\item
  Largest differences: cholesterol (+12.3) and triglycerides (+11.9)
\end{itemize}

\textbf{Box's M Test:}

\begin{itemize}
\item
  Corrected test fails to reject equal covariances (\(\chi^{2} = 8.46\),
  df = 10, p = 0.584)
\item
  No evidence against the MANOVA assumption -- not proof it holds
\end{itemize}
\end{frame}

\begin{frame}{Key Findings: Treatment Effectiveness}
\textbf{MANOVA Results (synthetic data):}

\begin{itemize}
\item
  Intervention improves all health outcomes (Λ = 0.889, p \textless{}
  0.0001)
\item
  Blood pressure: -6.5 / -4.2 mmHg
\item
  Cholesterol: -7.9 mg/dL
\item
  Glucose: -5.6 mg/dL
\end{itemize}

\textbf{Note:} Illustrates MANOVA mechanics on simulated data, not a
real treatment-effect finding
\end{frame}

\begin{frame}{Key Findings: Lifestyle-Health Relationships}
\textbf{Canonical Correlation:}

\begin{itemize}
\item
  Strong link between the first lifestyle and physiology variates (r =
  0.639)
\item
  Healthy lifestyle pattern: ↑ exercise, ↓ stress
\item
  Favorable health profile: ↓ BP, ↑ HDL
\item
  40.8\% shared variance between the \textbf{first canonical variate
  pair} -- not a domain-level ``lifestyle explains 41\% of health''
  statement, which would need variance-extracted/redundancy indices
\end{itemize}

\textbf{Insight (illustrative):} In this simulation, lifestyle and
health-marker patterns are linked as designed
\end{frame}

\begin{frame}{Methodological Insights}
\begin{enumerate}
\item
  \textbf{Multivariate methods reveal patterns} missed by univariate
  tests
\item
  \textbf{Type I error control} critical with multiple outcomes
\item
  \textbf{Assumption testing} (Box's M) validates results
\item
  \textbf{Effect sizes matter} beyond statistical significance
\item
  \textbf{Clinical context} guides interpretation
\end{enumerate}

All methods demonstrated with synthetic healthcare data generated for
this course
\end{frame}

\begin{frame}{Key Takeaways: Models}
\textbf{Logistic Regression:}

\begin{itemize}
\item
  Use for binary outcomes
\item
  Maximum likelihood estimation
\item
  Interpret via odds ratios
\item
  \textbf{Case Study:} 71\% accuracy predicting CVD risk
\end{itemize}
\end{frame}

\begin{frame}{Key Takeaways: Inference}
\textbf{Covariance Matrix Tests:}

\begin{itemize}
\item
  Box's M test for equality
\item
  Wishart distribution foundation
\item
  \textbf{Case Study:} M = 8.49 (assumption satisfied)
\end{itemize}

\textbf{Mean Vector Tests:}

\begin{itemize}
\item
  Hotelling's T-squared generalizes t-test
\item
  Confidence regions are ellipsoids
\item
  \textbf{Case Study:} T² = 228.65 (strong group differences)
\end{itemize}
\end{frame}

\begin{frame}{Key Takeaways: Advanced Methods}
\textbf{MANOVA:}

\begin{itemize}
\item
  Multiple response variables simultaneously
\item
  Wilks' Lambda most common test
\item
  Controls Type I error
\item
  \textbf{Case Study:} Λ = 0.889 (intervention effective)
\end{itemize}

\textbf{Canonical Correlation:}

\begin{itemize}
\item
  Relates two variable sets
\item
  Multiple correlation pairs
\item
  \textbf{Case Study:} r = 0.639 (lifestyle-health link)
\end{itemize}
\end{frame}

\begin{frame}{Key Takeaways: Applications}
\textbf{Factor-Based Regression:}

\begin{itemize}
\item
  Addresses multicollinearity
\item
  Dimension reduction
\item
  Interpretable factors
\end{itemize}

\textbf{Software:}

\begin{itemize}
\item
  Python: scikit-learn, statsmodels
\item
  R: stats, MASS
\item
  Commercial: SPSS, SAS
\item
  \textbf{Case Study:} All analyses implemented in Python
\end{itemize}
\end{frame}

\begin{frame}{Common Pitfalls to Avoid}
\begin{enumerate}
\item
  Using logistic regression without checking convergence
\item
  Ignoring multicollinearity in regression
\item
  Not checking MANOVA assumptions (Box's M)
\item
  Over-interpreting weak canonical correlations
\item
  Using too many factors in factor-based regression
\end{enumerate}
\end{frame}

\begin{frame}{Method Selection Guide}
{\def\LTcaptype{none} % do not increment counter
\begin{longtable}[]{@{}
  >{\raggedright\arraybackslash}p{(\linewidth - 2\tabcolsep) * \real{0.4286}}
  >{\raggedright\arraybackslash}p{(\linewidth - 2\tabcolsep) * \real{0.5714}}@{}}
\toprule\noalign{}
\endhead
\textbf{Situation} & \textbf{Method} \\
Binary outcome & Logistic regression \\
Multiple groups, multiple responses & MANOVA \\
Relate two variable sets & Canonical correlation \\
Multicollinear predictors & Factor/PCA regression \\
\bottomrule\noalign{}
\end{longtable}
}
\end{frame}

\begin{frame}{Recommended Resources: Books}
\textbf{Textbooks:}

\begin{itemize}
\item
  Agresti (2018) - Introduction to Categorical Data Analysis
\item
  Johnson \& Wichern (2007) - Applied Multivariate Statistical Analysis
\item
  Rencher \& Christensen (2012) - Methods of Multivariate Analysis
\end{itemize}
\end{frame}

\begin{frame}{Recommended Resources: Online}
\textbf{StatQuest YouTube Channel:}

\begin{enumerate}
\item
  \textbf{Logistic Regression:}\\
  \url{https://www.youtube.com/watch?v=yIYKR4sgzI8}
\item
  \textbf{MANOVA Concepts:} Search ``MANOVA StatQuest''
\item
  \textbf{PCA (for PCR):}\\
  \url{https://www.youtube.com/watch?v=FgakZw6K1QQ}
\end{enumerate}
\end{frame}

\begin{frame}{Recommended Resources: Software}
\textbf{Documentation:}

\begin{itemize}
\item
  Scikit-learn: \url{https://scikit-learn.org}
\item
  Statsmodels: \url{https://www.statsmodels.org}
\item
  R Documentation: \url{https://www.rdocumentation.org}
\end{itemize}
\end{frame}

\begin{frame}{Questions?}
\textbf{ Thank you for your attention! }

Juliho Castillo Colmenares\\
julihocc@tec

Office: Tec de Monterrey CCM Office 1540\\
Office Hours: Monday-Friday, 9:00 AM - 5:00 PM
\end{frame}

\begin{frame}{Next Steps: This Week}
\textbf{For This Week:}

\begin{itemize}
\item
  Review lecture notes thoroughly
\item
  Practice with provided examples
\item
  Complete practice questions
\item
  Prepare for E07 quiz
\end{itemize}
\end{frame}

\begin{frame}{Next Steps: Preparation for Evaluation}
\textbf{Key Topics to Master:}

\begin{itemize}
\item
  Logistic regression: logit transformation, MLE, interpretation
\item
  Hotelling's T-squared: computation and F transformation
\item
  MANOVA: assumptions, Wilks' Lambda, interpretation
\item
  Canonical correlation: number of pairs, loadings, significance
\item
  Factor-based regression: workflow, benefits, limitations
\item
  Software implementation: Python and R basics
\end{itemize}
\end{frame}

\begin{frame}{Integration with Previous Topics}
\textbf{Building on Earlier Concepts:}

\begin{itemize}
\item
  Factor Analysis (L04) → Factor-based regression
\item
  Discriminant Analysis (L05) → MANOVA post-hoc
\item
  PCA principles → Principal components regression
\end{itemize}

\textbf{Comprehensive Framework:} All methods part of the multivariate
toolkit
\end{frame}

\end{document}




